\problemname{Tunnelbaneplatser}

Baltic Olympiad in Informatics -- programmeringstävlingen för länderna i
östersjöregionen -- kommer till våren att anordnas i Stockholm, och
deltagarna kommer åka tunnelbana mellan vandrarhemmet och tävlingsarenan.

På tunnelbanetågen finns det sätesgrupper med fyra säten vardera. Nu
kommer ett antal grupper av deltagare och vill sätta sig. Varje grupp
har storlek 1, 2, 3 eller 4. Helst skulle alla personerna i en grupp
vilja sitta i samma fyrsätesgrupp, alltså slippa dela på sig. Hur många fyrsätesgrupper krävs för att detta ska vara möjligt?

\section*{Indata}

Indata består av en rad med fyra heltal $a_1, a_2, a_3, a_4$ -- antalet grupper av varje storlek. Alla tal är mellan $0$ och $100$.

\section*{Utdata}

Skriv ut ett enda tal: det minsta antalet fyrsätesgrupper som behövs för att
personerna ska kunna sätta sig så att alla inom varje grupp sitter tillsammans.

\section*{Förklaringar till exemplen}

\includegraphics[width=7cm]{tunnelbana1.pdf}\\
{\em En optimal placering i exempel 1 (streckade linjer avgränsar grupper av personer)}
\vspace{1cm}

\includegraphics[width=7cm]{tunnelbana2.pdf}\\
{\em En optimal placering i exempel 2.}
\vspace{1cm}

\includegraphics[width=7cm]{tunnelbana3.pdf}\\
{\em En optimal placering i exempel 3.}

\section*{Poängsättning}
För testfall värda $20$ poäng finns inga grupper av storlek 1.\\
För testfall värda $20$ poäng finns inga grupper av storlek 2.\\
För testfall värda $20$ poäng finns inga grupper av storlek 3.\\
För testfall värda $40$ poäng gäller inga speciella begränsningar.
